\documentclass[10pt,a4paper]{amsart}
\usepackage[margin=19mm]{geometry}
\usepackage{amsmath,amssymb,mathtools}
\usepackage[scr=rsfs,cal=euler]{mathalfa}
\usepackage{xfrac}
\usepackage[unicode,hidelinks,bookmarks=false]{hyperref}
\newcommand{\ie}{i.e.\ }
\newcommand{\cf}{cf.\ }
\newcommand{\resp}{respectively\ }
\newcommand{\Subsection}{\subsection}
\providecommand{\nobreakdash}{\nobreak\hbox{-}}
\setlength{\emergencystretch}{2em}
\multlinegap0pt
\usepackage{bm}
\usepackage{bbm}
\usepackage{dsfont}
\usepackage{xspace}
\usepackage{cancel}
\usepackage{mathtools}
\usepackage{cleveref}
\usepackage{amsthm}
\makeatletter
\@ifundefined{definition}{\newtheorem{definition}{Definition}}{}
\@ifundefined{lemma}{\newtheorem{lemma}[definition]{Lemma}}{}
\@ifundefined{proposition}{\newtheorem{proposition}[definition]{Proposition}}{}
\@ifundefined{remark}{\newtheorem{remark}[definition]{Remark}}{}
\makeatother
\newcommand{\BrTwo}{\mathrm{Br}_2}
\newcommand{\D}{\mathsf{D}}
\newcommand{\GP}{\mathsf{GP}}
\newcommand{\bipf}[1]{\langle\!\langle #1\rangle\!\rangle_b}
\newcommand{\dlog}[1]{\frac{\mathrm{dlog}#1}{#1}}
\newcommand{\inv}{\mathsf{inv}}
\newcommand{\norm}{\mathrm{N}}
\newcommand{\qpf}[1]{\langle\!\langle #1]\!]}
\newcommand{\quat}[1]{[#1)}
\title[A cohomological invariant for algebras of degree 8 and expon]{A cohomological invariant for algebras of degree 8 and exponent 2 in characteristic 2}
\author{Ahmed Laghribi, Nico Lorenz}
\date{Source: arXiv:2602.19675v1 (CC BY 4.0)}
\begin{document}
\maketitle

\noindent\emph{Source and licence.} A cohomological invariant for algebras of degree 8 and exponent 2 in characteristic 2. Version arXiv:2602.19675v1 (CC BY 4.0), distributed under the Creative Commons Attribution 4.0 International licence, as recorded in the arXiv licence metadata for this exact version and shown on the abstract page https://arxiv.org/abs/2602.19675v1. Full rights notice: licenses/arxiv-2602.19675-CC-BY-4.0.txt.\par\smallskip

\noindent\textit{Editorial scope.} Counted unit: the Lemma whose source label is \texttt{lem\_AlbertDefinedOverF} in the article (source0.tex lines 589--593, source label \texttt{lem\_AlbertDefinedOverF}), delivered together with the article's own complete original proof of it (source0.tex lines 594--630, one \texttt{proof} environment of 37 lines). No mathematics is written, repaired or re-derived: statement and proof are the article's own text line for line. Required context retained uncounted: prop\_DK\_ExistenceInv (lines 438--448); eq\_e3TraceSum (lines 531--536); rem\_EntriesOfBiquaternion (lines 553--566). Every definition, display and label of those lines is retained unchanged. Omitted from this package and not redistributed: 1--437, 449--530, 537--552, 567--588, 631--1244. Nothing inside a retained excerpt is omitted or altered: every retained body is delivered exactly as the article's lines are written, so no omission receipt of any kind accompanies this package. Citations: the retained excerpts carry no citation command at all, so no bibliography entry of the article is transcribed and no reference is redistributed. Reference adaptation: none was needed and none was made; every reference inside the delivered unit resolves inside it, so no unresolved reference or citation is produced. Printed slips kept literal: no printed word, symbol, hypothesis or number of the counted statement or of its proof was altered. Layout and class: the article's own class is \texttt{amsart} with packages including inputenc, amsthm, amsmath, amssymb, enumitem, tikz, xcolor, mydiagrams, todonotes, hyperref, cleveref; the wrapper is the lane's pinned \texttt{amsart} editorial wrapper at 10pt on A4, which restates the theorem-like environments the retained text needs and adds the article's own macro definitions that the retained text needs (\texttt{\textbackslash BrTwo}, \texttt{\textbackslash D}, \texttt{\textbackslash GP}, \texttt{\textbackslash bipf}, \texttt{\textbackslash dlog}, \texttt{\textbackslash inv}, \texttt{\textbackslash norm}, \texttt{\textbackslash qpf}, \texttt{\textbackslash quat}), with extra wrapper packages (bm, bbm, dsfont, xspace, cancel, mathtools, cleveref). The delivered numbering is the unit's own. Assets: no retained span contains a figure environment, a graphics-inclusion command, a PDF-page inclusion, a table or a TikZ picture; no image file is copied into this package, the package declares no asset dependency and no image file is redistributed with it. Licence: version arXiv:2602.19675v1 of this article is distributed under the Creative Commons Attribution 4.0 International licence, as recorded in the arXiv licence metadata for this exact version and shown on the abstract page https://arxiv.org/abs/2602.19675v1. A full rights notice is delivered at licenses/arxiv-2602.19675-CC-BY-4.0.txt. Characters: every retained excerpt is pure ASCII, so the delivered engine, XeLaTeX with its default Latin Modern text font, reports no missing character.
\begin{proposition} \label{prop_DK_ExistenceInv}
    Let $[D] \in \BrTwo(F)$ be such that $[D_L] = 0$.

    Then, there are $x, y \in K^\ast$ such that
    \[[D_K] = [\quat{a, x} \otimes \quat{b, y}]\]
    and 
    \begin{align*}
        \inv(D) &:= e^3(s_\ast(\qpf{x, a} \perp \qpf{y, b})) \in H_2^3(F) / (N(d; a, b) \wedge [D]),
    \end{align*}
    does not depend on the choice of $x, y$.
\end{proposition}

    \begin{align} \label{eq_e3TraceSum}
        e^3(s_\ast(\varphi)) + e^3(s_\ast(\varphi')) &= e^3(s_\ast(\varphi \perp \varphi'))\nonumber \\
        &= s_\ast(e^3(\varphi \perp \varphi')) \nonumber\\
        &= s_\ast(e^3(\bipf{\norm_{E / K}(\gamma)} \otimes \varphi)) \nonumber\\
        &= s_\ast\left({\overline{\dlog\norm_{E / K}(\gamma)}} \wedge [D_K] \right).
    \end{align}

\begin{remark} \label{rem_EntriesOfBiquaternion}
	The argument in \Cref{prop_DK_ExistenceInv} shows, that if $\gamma \in E^\ast$ occuring in \Cref{eq_e3TraceSum} fulfils $\qpf{xx', a} \cong \qpf{\norm_{E/K}(\gamma), a}$, then 
	\[e_3(s_\ast(\varphi)) + e_3(s_\ast(\varphi')) = \overline{\dlog(\norm_{E/F}(\gamma))} \wedge [D].\]
   % \begin{enumerate}
%        \item The existence of $x, y \in K^\ast$ in \Cref{eq_D_KAlbertForm} also follows from a deep theorem of R. Aravire and B. Jacob without using an Albert form, see \cite[Theorem 20]{AravireJacob_GradedKernelBiquadraticSeparable}.
%        \item By comparing the Arf invariants 
%        \[0 = \Delta(\varphi) = \Delta(\rho_1\otimes[1, a] \perp \rho_2\otimes [1, b]) = \dim(\rho_1)\cdot a + \dim(\rho_2) \cdot b\] 
%        and using $[L : F] = 8$, we readily see 
%        \[\dim(\rho_1) \equiv \dim(\rho_2) \equiv 0 \mod 2.\]
%        Thus, one can calculate explicitly $x = \det(\rho_1)$ and $y = \det(\rho_2)$.
       % \item\label{rem_EntriesOfBiquaternion3} The argument in \Cref{prop_DK_ExistenceInv} shows, that the element $\gamma$ occuring in \Cref{eq_e3TraceSum} is determined by the isometry
       % \[\qpf{xx', a} \cong \qpf{\norm_{E/K}(\gamma), a}.\]
    %\end{enumerate} 
\end{remark}

\begin{lemma} \label{lem_AlbertDefinedOverF}
    If $\inv(D) = 0$, then there are $x, y \in K^\ast$ such that $[D_K] = [\quat{a, x} \otimes \quat{b, y}$ and the corresponding Albert form
   \[[1, a + b] \perp x [1, a] \perp y[1, b]\]
    is defined over $F$. 
\end{lemma}

\begin{proof}
    Suppose $\inv(D) = 0$.
    We have to prove that $[D_K] = [\quat{a, x} \otimes \quat{b, y}$ for some $x, y \in K^\ast$ such that the corresponding Albert form is defined over $F$.
    By the first part of \Cref{prop_DK_ExistenceInv}, there are $\Tilde x, \Tilde y \in K^\ast$ such that 
    \[[D_K] = [\quat{a, \Tilde x} \otimes \quat{b, \Tilde y}].\]
    We consider the Albert form
    \[\Tilde\varphi = [1, a + b] \perp \Tilde x [1, a] \perp \Tilde y[1, b]\]
    corresponding to $\quat{a, \Tilde x} \otimes \quat{b, \Tilde y}$.
    If $\Tilde\varphi$ is defined over $F$, we are done.
    
    Otherwise, since $\inv(D) = 0$, by definition of $\inv$, there is $\gamma \in E^\ast$ such that 
    \begin{align} \label{eq_ChoiceGamma}
        e^3(s_\ast(\tilde \varphi)) = \overline{\dlog(\norm_{E/F}(\gamma))} \wedge [D]).
    \end{align}
    Let $x = \Tilde x^{-1} \cdot \norm_{E/K}(\gamma)$ and $ y = \Tilde y^{-1} \cdot \norm_{E/K}(\gamma)$.
    Using
    \begin{align}\label{eq_normRepresented}
        \norm_{E/K}(\gamma) \in \D_K([1, a + b]),
    \end{align}
    in $\BrTwo(K)$ we have 
    \begin{align*}
        \quat{a, x} \otimes \quat{b, y} \otimes \quat{a, \Tilde{x}} \otimes \quat{b, \Tilde{y}} &= \quat{a, x\Tilde{x}} \otimes \quat{b, y\Tilde{y}}\\
        &= \quat{a, \norm_{E/F}(\gamma)} \otimes \quat{b, \norm_{E/K}(\gamma)}\\
        &= \quat{a + b, \norm_{E/K}(\gamma)} = 0
    \end{align*}
    and thus 
    \[\quat{a, x} \otimes \quat{b, y} \cong \quat{a, \Tilde{x}} \otimes \quat{b, \Tilde{y}}.\]
    Therefore, 
    \[{\varphi} = [1, a + b] \perp x [1, a] \perp y[1, b]\]
    is also an Albert form of $D_K$.
    Since we have $\quat{a, x\Tilde{x}} \cong \quat{a, \norm_{E/K}(\gamma)}$, we have 
    \[e^3(s_\ast(\varphi)) + e^3(s_\ast(\Tilde{\varphi})) = \overline{\dlog \norm_{E/F}(\gamma)} \wedge [D],\]
    by \Cref{rem_EntriesOfBiquaternion} and by the choice of $\gamma$.
    Using \eqref{eq_ChoiceGamma} we obtain 
    \[e^3(s_\ast({\varphi})) = 0.\]
    Now since $s_\ast({\varphi}) \in \GP_3(F)$, this form is hyperbolic by the Arason-Pfister Hauptsatz, so ${\varphi}$ is defined over $F$.
\end{proof}

\end{document}
